\appendix
\section{Detailed Experimental Setup}
\subsection{WeaveBench}
WeaveBench~\cite{weavebench} consists of 114 long-horizon computer-use tasks spanning eight domains: Desktop Applications, Document Processing, Games, Web Development, Data Analysis and Visualization, DevOps, Spatial/3D Applications, and Design. Each task requires the agent to use both GUI and CLI interactions within the same workflow, rather than relying solely on command-line or browser operations. We follow the standard WeaveBench evaluation protocol: each task is launched in an isolated containerized Ubuntu desktop VM, and task executions do not share runtime state. The task definitions, workspaces, runtime assets, judge templates, and scoring procedure are all taken from the official release. The evaluator uses Claude Opus 4.7 as a trajectory-aware judge. We report two aggregate metrics: PassRate, defined as the fraction of tasks with score at least 0.8, and Overall, defined as the average score over all 114 tasks. In addition to the aggregate results, we report domain-level PassRate for the eight domains: Desktop Applications (DSK), Document Processing (DOC), Games (GAM), Web Development (WEB), Data Analysis and Visualization (DAV), DevOps (OPS), Spatial/3D Applications (SPA), and Design (DES).

For the model and execution framework, we evaluate Qwen 3.7-Plus. The Claude Code runtime version is \texttt{claude-code 2.1.76}. Each task is allowed up to 25 Manage-Execute-Audit rounds. The per-round timeout is 1800 seconds for the executor, and 300 seconds for both the manager and the verifier. The Claude Code reasoning effort on the task side is set to \texttt{high}; the judge is executed with OpenClaw, with \texttt{AJ\_THINKING=medium}.

For GUI tools, we preserve the basic WeaveBench computer-use interface, including the \texttt{screenshot} perception tool and desktop actuation tools implemented with \texttt{pyautogui}, such as click, double-click, mouse move, keyboard input, scroll, drag, and wait. In addition, we introduce a restricted evidence-preservation tool, \texttt{save\_screenshot}. This tool does not provide any new environment-control capability, nor does it expose additional information to the agent; it can only save the currently visible real desktop screen to the target directory as an official screenshot artifact required by certain tasks.

Since GUI tasks frequently provide screenshot inputs, we use an image URL proxy for visual inputs: screenshots are uploaded to an image-hosting backend and referenced as image URLs in model requests, avoiding oversized request bodies caused by multiple base64-encoded screenshots.

\subsection{OSWorld 2.0}
OSWorld 2.0~\cite{osworld2} consists of 108 professional desktop workflow tasks, with a median human completion time of approximately 1.6 hours. We use the official OSWorld-v2 \texttt{osworld-v2-2026.06.24} release, including the corresponding task definitions, task assets, Docker VM image, and mocked websites. Experiments are conducted using the standard OSWorld Docker-based VM infrastructure in an Ubuntu desktop environment, with a default resolution of $1920\times1080$. Final scoring is performed by the native OSWorld \texttt{env.evaluate()} procedure.

Unlike the main official baselines, which typically rely only on GUI actions, LongHorizon-Harness uses a hybrid tool pool on OSWorld-v2. GUI actions are executed by a computer MCP server inside the VM, supporting desktop operations such as clicking, typing, dragging, and hotkeys. The CLI role can execute shell commands in the same VM for filesystem inspection, script execution, and long-context task processing. These tools are hosted inside the VM through a fixed version of \texttt{claude-code 2.1.176}.

For model inference, we set \texttt{thinking.type=enabled}, \texttt{max\_tokens}=65{,}536, and \texttt{effort=max}. The native OSWorld evaluator and user simulator use the same \texttt{qwen3.7-plus} configuration as the execution agent. For human-in-the-loop tasks in the benchmark that require user information or authorization, we do not introduce manual intervention; instead, we use the native OSWorld user-simulator interface to provide responses.

We report two metrics: Binary Accuracy and Partial Accuracy. Binary Accuracy counts a task as successful only when its final score is 1, while Partial Accuracy is the average fine-grained score over all 108 tasks and measures partial task completion.

\subsection{Terminal-Bench 2.1}
Terminal-Bench 2.1~\cite{terminalbench} is designed to evaluate agents on challenging and realistic software-engineering and command-line tasks. We use Harbor as the evaluation framework, with the Docker backend providing an isolated execution environment for each task. During evaluation, we preserve the CPU, memory, and environment constraints originally defined for each task.

To ensure consistency with existing evaluations of Qwen-series models, both the baseline and CUA-Harness settings use Claude Code as the underlying CLI execution agent, connected to the \texttt{qwen3.7-plus} model. The Claude Code version is \texttt{v2.1.211}. Since the Claude Code command-line interface does not directly expose all low-level request parameters, we introduce a transparent request proxy inside the container. All model requests are configured with \texttt{temperature}=1.0, \texttt{top\_p}=0.95, \texttt{top\_k}=20, \texttt{enable\_thinking=true}, and \texttt{max\_tokens}=65{,}536.

Each trial has a timeout of 5 hours. We run three independent trials per task, report the average score over the three runs for each task, and compare against the Claude Code baseline under the same model setting as well as the official leaderboard results.

\section{Detailed Experimental Results}
\subsection{Detailed Results on the OSWorld 2.0 Opus 4.7 Subset}

Table~\ref{tab:osworld-opus47-detailed} reports the per-task scores on the 34-task OSWorld 2.0 Opus 4.7 subset. The baseline uses the standard single-action GUI setting, while LongHorizon-Harness uses our hybrid GUI+CLI tool pool. On this subset, LongHorizon-Harness improves the partial score from $55.83\%$ to $66.86\%$, and the binary accuracy from $20.6\%$ to $35.3\%$.

\begin{table*}[t]
\centering
\caption{Detailed per-task results on the OSWorld 2.0 Opus 4.7 subset. Ours denotes LongHorizon-Harness.}
\setlength{\tabcolsep}{7pt}
\begin{tabular}{lrrlrr}
\toprule
Task & Baseline & Ours & Task & Baseline & Ours \\
\midrule
001 & 0.7778 & 1.0000 & 018 & 0.0000 & 0.0000 \\
002 & 0.7500 & 1.0000 & 019 & 0.5200 & 0.7600 \\
003 & 1.0000 & 1.0000 & 020 & 0.1000 & 1.0000 \\
004 & 0.8000 & 0.7667 & 021 & 1.0000 & 0.0000 \\
005 & 0.0000 & 1.0000 & 022 & 1.0000 & 1.0000 \\
006 & 0.5556 & 0.8889 & 023 & 0.9300 & 0.9300 \\
007 & 0.4000 & 0.8000 & 025 & 0.2401 & 0.0833 \\
010 & 1.0000 & 1.0000 & 027 & 0.5000 & 0.0000 \\
011 & 0.2963 & 0.5556 & 028 & 0.4000 & 1.0000 \\
012 & 0.6000 & 1.0000 & 030 & 1.0000 & 0.0000 \\
013 & 1.0000 & 0.7500 & 032 & 0.9111 & 0.2000 \\
014 & 0.9000 & 0.7400 & 033 & 0.1000 & 1.0000 \\
015 & 0.2840 & 0.7667 & 034 & 0.5938 & 0.8750 \\
016 & 0.8571 & 0.7524 & 051 & 0.0000 & 0.6000 \\
017 & 0.1429 & 0.0000 & 052 & 1.0000 & 1.0000 \\
060 & 0.1000 & 0.1000 & 054 & 0.3628 & 1.0000 \\
061 & 0.3574 & 0.3891 & 062 & 0.5033 & 0.7753 \\
\bottomrule
\end{tabular}

\label{tab:osworld-opus47-detailed}
\end{table*}

\subsection{Fine-Grained Breakdown by Benchmark and Task Type}
\label{app:fine-grained-breakdown}

Beyond the Opus 4.7 subset reported above, we provide a fine-grained breakdown of the main evaluation results by benchmark, domain, task property, and failure mode. We use LongHorizon-Harness (LH-Harness) as the abbreviated name in the tables. The purpose of this analysis is to identify where the harness changes the execution dynamics, rather than only reporting aggregate gains. Across benchmarks, the largest improvements appear when progress can be represented as verifiable environment state: installed binaries, repository states, screenshots with well-defined metadata, application configuration, structured files, or cross-checked evidence from multiple interfaces. The remaining failures concentrate in tasks where the decisive condition is difficult to close under the available verifier, such as hidden performance thresholds, embodied visual precision, temporal video evidence, or task semantics that admit several plausible interpretations.

\subsubsection{WeaveBench}
\label{app:fine-grained-weavebench}

Table~\ref{tab:weavebench-domain-breakdown} reports the domain-level breakdown on WeaveBench. LH-Harness improves the mean score in seven of the eight domains, with the largest gains in Design (DES), Spatial/3D Applications (SPA), and Games (GAM). Pass rate also increases in every domain, including Desktop (DSK), where the mean score slightly decreases despite a higher pass rate. This pattern indicates that the harness mainly raises the lower tail of long-horizon trajectories: it converts many near-failures into partial or complete successes, while a small number of already-strong baseline trajectories can be affected by additional verification and repair steps.

\begin{table}[t]
\caption{WeaveBench domain-level results. Score $\Delta$ is LH-Harness minus baseline. Pass-rate improvement is reported in percentage points.}
\label{tab:weavebench-domain-breakdown}
\centering
\small
\setlength{\tabcolsep}{5pt}
\begin{tabular}{lrrrr}
\toprule
Domain & Baseline & LH-Harness & Score $\Delta$ & Pass-rate $\Delta$ \\
\midrule
DES & 0.5630 & 0.8160 & +0.2530 & +60.00 pp \\
SPA & 0.5800 & 0.8192 & +0.2392 & +50.00 pp \\
GAM & 0.5235 & 0.7328 & +0.2093 & +29.41 pp \\
DAV & 0.6981 & 0.8608 & +0.1627 & +30.77 pp \\
OPS & 0.7758 & 0.9090 & +0.1332 & +25.00 pp \\
DOC & 0.8009 & 0.8929 & +0.0920 & +23.53 pp \\
WEB & 0.7287 & 0.8147 & +0.0860 & +26.67 pp \\
DSK & 0.8671 & 0.8465 & -0.0205 & +5.56 pp \\
\bottomrule
\end{tabular}
\end{table}

The Games domain provides a useful stress test because it combines visual interaction, symbolic state, search, and environment-specific constraints. Table~\ref{tab:weavebench-game-breakdown} shows that the largest gains occur on tasks where the baseline often fails to establish a stable state representation. For example, Mines, Stockfish puzzle analysis, Quadrapassel autoplay, PokerTH equity play, and SuperTux level repair all improve from near-zero or zero baseline scores to nontrivial completion. In these cases, the harness helps by decomposing the task into explicit subgoals and by carrying forward verified facts rather than a long interaction transcript. In contrast, when the baseline already solves the core task, additional rounds provide less benefit and can even slightly reduce score, as in rhythm autoplay, X-Moto, KMahjongg, and FreeCell.

\begin{table*}[t]
\caption{Fine-grained results for WeaveBench Games tasks. Task names omit the shared \texttt{GAM\_task\_*} prefix for compactness.}
\label{tab:weavebench-game-breakdown}
\centering
\scriptsize
\setlength{\tabcolsep}{4pt}
\begin{tabular}{lrr@{\hspace{1.2em}}lrr}
\toprule
Task & Baseline & LH-Harness & Task & Baseline & LH-Harness \\
\midrule
gnome\_mines\_solve & 0.040 & 0.590 & rhythm\_autoplay & 0.950 & 0.830 \\
game\_ui\_bug & 0.750 & 0.850 & gnuchess\_pgn\_blunder\_hunt & 0.830 & 0.940 \\
stockfish\_puzzle\_analysis & 0.000 & 0.550 & anagramarama\_word\_grid & 0.750 & 0.830 \\
gdb\_pygame\_cheat & 0.720 & 0.860 & godot\_scene\_node\_debug & 0.930 & 0.870 \\
mines\_visual & 0.000 & 0.470 & hedgewars\_lua\_mission\_debug & 0.740 & 0.830 \\
quadrapassel\_autoplay & 0.000 & 0.300 & supertux\_level\_repair\_play & 0.000 & 0.650 \\
pokerth\_equity\_play & 0.000 & 0.920 & xmoto\_level\_xml\_repair\_ride & 0.920 & 0.850 \\
sokoban\_solver & 0.960 & 0.958 & kmahjongg\_pair\_solver & 0.790 & 0.750 \\
 & & & pysol\_freecell\_fcsolve & 0.520 & 0.410 \\
\bottomrule
\end{tabular}
\end{table*}

\subsubsection{OSWorld 2.0}
\label{app:fine-grained-osworld}

Table~\ref{tab:osworld-tag-breakdown} reports the OSWorld 2.0 breakdown by capability tag. Tags are not mutually exclusive, so a task may contribute to multiple rows. LH-Harness improves all six tag groups, but the magnitude and residual failure modes differ substantially. The largest gain appears in streaming interaction tasks, where the baseline score is zero and LH-Harness reaches 0.500 on average. Human-in-the-loop tasks also improve strongly, indicating that the harness benefits from converting user-provided information into persistent state constraints. Tutorial-following tasks improve because symbolic instructions can often be decomposed into bounded subtasks with explicit acceptance criteria.

\begin{table}[t]
\caption{OSWorld 2.0 results by capability tag. Tags are overlapping. ``Zero'' and ``Full'' count tasks in the tag group where LH-Harness obtains score 0 or 1, respectively.}
\label{tab:osworld-tag-breakdown}
\centering
\small
\setlength{\tabcolsep}{4pt}
\begin{tabular}{lrrrrrr}
\toprule
Tag & $n$ & Baseline & LH-Harness & $\Delta$ & Zero & Full \\
\midrule
streaming\_interaction & 6 & 0.000 & 0.500 & +0.500 & 3 & 3 \\
human\_in\_the\_loop & 6 & 0.225 & 0.557 & +0.332 & 1 & 0 \\
tutorial\_following & 22 & 0.157 & 0.374 & +0.217 & 7 & 2 \\
implicit\_state\_inference & 43 & 0.241 & 0.366 & +0.125 & 15 & 6 \\
visual\_spatial\_precision & 45 & 0.198 & 0.298 & +0.100 & 14 & 4 \\
cross\_source\_reasoning & 46 & 0.263 & 0.361 & +0.098 & 11 & 3 \\
\bottomrule
\end{tabular}
\end{table}

The tag-level results suggest five more specific observations. First, in human-in-the-loop tasks, the main benefit comes from treating the user's response as a durable constraint rather than as a transient dialogue turn. The same mechanism can also fail: if newly supplied information does not override stale evidence, the verifier may confirm an outdated state. Second, the harness is stronger at following symbolic tutorials than embodied tutorials. Commands, tables, file names, and web fields can be decomposed into auditable steps; video timelines, CAD geometry, and fine mouse operations remain difficult to translate into reliable low-level actions. Third, implicit state inference works best when there is an authoritative closure over the environment state, such as a file, database entry, application setting, or evaluator-visible artifact. Without such closure, plausible evidence can be mistaken for completed evidence. Fourth, visual-task gains mainly arise when a visual requirement can be rewritten as a file, code, metadata, or search problem. When the requirement remains purely spatial or motor-level, performance is still limited by visual localization and cursor control. Fifth, cross-source reasoning is often not the bottleneck by itself. Many failures occur after the relevant evidence has been understood, when the final answer is written to the wrong state carrier, saved at the wrong boundary, or not submitted through the expected interface.

\subsubsection{Terminal-Bench 2.1}
\label{app:fine-grained-terminalbench}

Terminal-Bench provides a complementary view because most tasks are command-line tasks with hidden tests. Table~\ref{tab:terminalbench-category-breakdown} reports the category-level results. The largest gain is in system administration, where LH-Harness improves from 0.593 to 0.889. This category is dominated by stateful procedures such as installation, service configuration, PATH management, build outputs, and persistent environment side effects. These are precisely the settings in which an explicit task contract and independent audit can prevent the agent from stopping after an incomplete but superficially plausible command sequence. Software engineering also improves substantially, especially in tasks where the baseline performs most of the work but misses one blocking constraint, such as image similarity, asynchronous cancellation behavior, independence from a reference binary, or compressed-size limits.

\begin{table*}[t]
\caption{Terminal-Bench 2.1 category-level results.}
\label{tab:terminalbench-category-breakdown}
\centering
\small
\setlength{\tabcolsep}{5pt}
\begin{tabular}{lrrrr}
\toprule
Category & $n$ & Baseline & LH-Harness & $\Delta$ \\
\midrule
software-engineering & 78 & 0.705 & 0.833 & +0.128 \\
system-administration & 27 & 0.593 & 0.889 & +0.296 \\
data-science & 24 & 0.792 & 0.667 & -0.125 \\
scientific-computing & 24 & 0.375 & 0.542 & +0.167 \\
security & 24 & 0.833 & 0.875 & +0.042 \\
file-operations & 15 & 0.333 & 0.333 & +0.000 \\
debugging & 15 & 0.933 & 1.000 & +0.067 \\
data-processing & 12 & 0.750 & 0.917 & +0.167 \\
mathematics & 12 & 0.917 & 0.750 & -0.167 \\
model-training & 12 & 0.750 & 0.667 & -0.083 \\
machine-learning & 9 & 1.000 & 1.000 & +0.000 \\
games & 3 & 0.000 & 0.333 & +0.333 \\
video-processing & 3 & 0.333 & 0.000 & -0.333 \\
data-querying & 3 & 1.000 & 1.000 & +0.000 \\
optimization & 3 & 1.000 & 1.000 & +0.000 \\
personal-assistant & 3 & 1.000 & 1.000 & +0.000 \\
\bottomrule
\end{tabular}
\end{table*}

Table~\ref{tab:terminalbench-difficulty-breakdown} shows that the gain is larger on hard tasks than on medium tasks. This is consistent with the purpose of the harness: when a task is short enough for the baseline to complete in one trajectory, the marginal value of explicit state management is smaller. As tasks become longer and contain more intermediate failure points, the ability to record unresolved constraints and re-plan from verified state becomes more important.

\begin{table}[t]
\caption{Terminal-Bench 2.1 results by difficulty.}
\label{tab:terminalbench-difficulty-breakdown}
\centering
\small
\setlength{\tabcolsep}{7pt}
\begin{tabular}{lrrrr}
\toprule
Difficulty & $n$ & Baseline & LH-Harness & $\Delta$ \\
\midrule
easy & 12 & 0.833 & 1.000 & +0.167 \\
medium & 165 & 0.776 & 0.818 & +0.042 \\
hard & 90 & 0.533 & 0.656 & +0.122 \\
\bottomrule
\end{tabular}
\end{table}

The tag-level results in Table~\ref{tab:terminalbench-tag-breakdown} further localize the effect. The strongest positive tags are images, version-control, system, and sys-admin. These tags share a common property: correctness can be checked against durable artifacts, such as image hashes or similarity scores, repository history, installed executables, service state, logs, or filesystem layout. Negative tags such as mteb, data-science, and video-processing expose a different limitation. In these tasks, the official correctness condition can depend on hidden thresholds, ranking semantics, temporal localization, or evaluation conventions that are not fully recoverable from visible state. In such cases, LH-Harness may organize execution more carefully, but a misinterpreted contract can still lead to a confidently verified wrong answer.

\begin{table*}[t]
\caption{Terminal-Bench 2.1 tag-level results. Tags may overlap across tasks.}
\label{tab:terminalbench-tag-breakdown}
\centering
\scriptsize
\setlength{\tabcolsep}{4pt}
\begin{tabular}{lrrrr@{\hspace{1.1em}}lrrrr}
\toprule
Tag & $n$ & Baseline & LH-Harness & $\Delta$ &
Tag & $n$ & Baseline & LH-Harness & $\Delta$ \\
\midrule
coding & 54 & 0.907 & 0.926 & +0.019 & sys-admin & 9 & 0.667 & 0.889 & +0.222 \\
file-operations & 30 & 0.733 & 0.767 & +0.033 & web & 9 & 0.556 & 0.667 & +0.111 \\
system & 30 & 0.633 & 0.900 & +0.267 & images & 6 & 0.167 & 1.000 & +0.833 \\
security & 27 & 0.815 & 0.778 & -0.037 & mteb & 6 & 0.500 & 0.167 & -0.333 \\
data-processing & 21 & 0.762 & 0.714 & -0.048 & video-processing & 6 & 0.333 & 0.000 & -0.333 \\
software-engineering & 21 & 0.714 & 0.714 & +0.000 & image-processing & 6 & 0.667 & 0.833 & +0.167 \\
compilation & 12 & 1.000 & 0.917 & -0.083 & ocr & 6 & 0.667 & 0.833 & +0.167 \\
machine-learning & 12 & 1.000 & 0.917 & -0.083 & log-analysis & 6 & 0.833 & 1.000 & +0.167 \\
version-control & 12 & 0.667 & 1.000 & +0.333 & simulation & 6 & 0.833 & 1.000 & +0.167 \\
biology & 9 & 0.000 & 0.111 & +0.111 & corewars & 6 & 1.000 & 0.833 & -0.167 \\
cloning & 9 & 0.000 & 0.111 & +0.111 & gaming & 6 & 1.000 & 0.833 & -0.167 \\
data-science & 9 & 0.667 & 0.333 & -0.333 & pmars & 6 & 1.000 & 0.833 & -0.167 \\
 & & & & & pytorch & 6 & 1.000 & 0.833 & -0.167 \\
\bottomrule
\end{tabular}
\end{table*}



\section{Additional Case Studies}
\label{app:additional-case-studies}

This appendix expands the qualitative analysis in Section~3.4 with additional case studies from WeaveBench and Terminal-Bench. These examples are selected to illustrate the same mechanism that underlies the quantitative gains: LongHorizon-Harness does not merely give the model more actions, but restructures long-horizon execution around explicit task state, fresh-context execution, and independent verification. Across domains, the repeated pattern is that local execution progress becomes useful only after it is converted into an auditable state proposition: what changed in the environment, which requirement it satisfies, what evidence supports it, and what remains unresolved.

\paragraph{What the cases are meant to show.}
The cases are not intended as isolated demonstrations of tool use. Instead, they expose recurring failure modes in long-horizon agents and show how the Manage-Execute-Audit loop addresses them. First, GUI failures are externalized as unresolved state rather than being trapped inside a growing interaction history. Second, artifacts are not accepted because the executor claims completion; they are accepted only after a read-only auditor inspects the resulting environment. Third, hybrid GUI+CLI execution lets the agent use the interface that best grounds each part of the task: GUI tools for visual and application state, and CLI tools for files, logs, metadata, build products, and reproducibility checks. Finally, the manager turns the audit result into the next bounded subtask, preserving verified progress while avoiding context rot.

\subsection{Cross-Domain WeaveBench Cases}
\label{app:weavebench-additional-cases}

\paragraph{Desktop workflow: turning a GUI failure into a verified state transition.}
Figure~\ref{fig:case-dsk} shows a desktop application task in Joplin. The task requires creating a rich note, placing it in the correct notebook, and reproducing a title-bar rendering issue. The important point is not that the executor eventually finds a way to click through the interface. The important point is that failed shortcuts and imprecise clicks are treated as changes in the task state: the note body already exists, the notebook still needs to be created, and the final visual bug still needs evidence. LongHorizon-Harness preserves the verified intermediate artifact, routes the remaining GUI work as a focused subtask, and then uses pixel-level screenshot comparison to audit the final rendering difference. The case illustrates failure recovery without discarding completed work or trusting the executor's self-assessment.

\begin{figure*}[t]
  \centering
  \includegraphics[width=\textwidth]{pic/DSK_en.pdf}
  \caption{Desktop workflow case. LongHorizon-Harness creates the rich note, recovers from failed GUI interactions by updating the task state, and verifies the reproduced rendering bug through screenshot evidence.}
  \label{fig:case-dsk}
\end{figure*}

\paragraph{Document processing: verifying semantic document state rather than visual appearance.}
Figure~\ref{fig:case-doc} illustrates a document-editing task where the visible appearance of headings is insufficient evidence of completion. A heading can look bold and large while still using the wrong underlying style. LongHorizon-Harness first inspects the real LibreOffice Writer state, then applies a reproducible macro to normalize headings, and finally audits the resulting ODT XML. This turns a visually plausible edit into an independently checkable proposition: all required headings are encoded with the correct heading style and outline level. The case highlights why completion must be grounded in environment state rather than screenshots or executor confidence alone.

\begin{figure*}[t]
  \centering
  \includegraphics[width=\textwidth]{pic/DOC_en.pdf}
  \caption{Document-processing case. The harness verifies the underlying ODT structure instead of relying on visual formatting, ensuring that all headings are semantically normalized.}
  \label{fig:case-doc}
\end{figure*}

\paragraph{Game analysis: combining GUI replay with file-level validation.}
Figure~\ref{fig:case-gam} shows a chess-analysis task in XBoard. The agent must identify an illegal move and analyze the game state around the failure. LongHorizon-Harness uses the GUI to replay the PGN until XBoard itself stops at the illegal move, preserving visual evidence of the exact board position. It then complements this GUI evidence with file-level checks of the generated reports. The result is a cross-validated artifact: the illegal move, the stopped board state, the blunder analysis, and the output files agree with one another. This case demonstrates the value of hybrid execution when neither GUI observation nor file inspection alone is sufficient.

\begin{figure*}[t]
  \centering
  \includegraphics[width=\textwidth]{pic/GAM_en.pdf}
  \caption{Game-analysis case. GUI replay identifies the illegal PGN transition, while file-level auditing checks that the final reports are consistent with the board evidence.}
  \label{fig:case-gam}
\end{figure*}

\paragraph{Web diagnostics: converging browser charts, JSON state, and packet evidence.}
Figure~\ref{fig:case-web} shows a WebRTC simulcast-layer audit. The task requires identifying a layer degradation event and supporting it with evidence from the browser and packet capture. LongHorizon-Harness separates the workflow into evidence-producing subtasks: render the browser charts, hover at the relevant timestamp, capture tooltip values, and open Wireshark to cross-check the RTP streams. The auditor then treats each screenshot as one proposition in a larger evidence chain, such as the f-layer dropping to zero frames per second around the target interval. This case illustrates how the harness prevents a long GUI investigation from becoming an unstructured transcript; each observation is recorded as a verified fact that can guide the next step.

\begin{figure*}[t]
  \centering
  \includegraphics[width=\textwidth]{pic/WEB_en.pdf}
  \caption{Web-diagnostics case. LongHorizon-Harness combines browser charts, tooltip evidence, and Wireshark inspection into a single audited evidence chain.}
  \label{fig:case-web}
\end{figure*}

\paragraph{Data analysis and visualization: auditing evidence quality, not just file existence.}
Figure~\ref{fig:case-dav} shows an Airflow debugging task. A naive trajectory can save screenshots with the right filenames while capturing the wrong UI state, such as a Grid view saved under a Gantt-view name. LongHorizon-Harness detects this mismatch during audit, keeps the noncompliant screenshot from supporting completion, and assigns a focused repair subtask to capture the real Gantt view. It also cross-checks the DAG graph to identify the incorrect upstream dependency. The key lesson is that artifact presence is weaker than artifact validity: a file exists only becomes progress after the auditor confirms that it depicts the required state.

\begin{figure*}[t]
  \centering
  \includegraphics[width=\textwidth]{pic/DAV_en.pdf}
  \caption{Data-analysis case. The harness rejects mislabeled evidence, captures the correct Airflow views, and audits whether each screenshot supports the required diagnosis.}
  \label{fig:case-dav}
\end{figure*}

\paragraph{DevOps: linking symptoms, root cause, and repair across UI and CLI evidence.}
Figure~\ref{fig:case-ops} presents a RabbitMQ operations task. The surface symptom is visible in the management UI: messages are published faster than they are delivered or acknowledged. LongHorizon-Harness then combines UI inspection with queue and binding evidence, identifying that the dead-letter routing key is mismatched. After repair, the auditor cross-validates the queue depths and binding configuration against CLI dumps and management UI screenshots. This is a representative operations workflow: the answer is not a single command, but a chain from symptom, to root cause, to repair, to post-repair evidence.

\begin{figure*}[t]
  \centering
  \includegraphics[width=\textwidth]{pic/OPS_en.pdf}
  \caption{DevOps case. The harness links management-UI symptoms to queue and binding evidence, then verifies the repaired routing state across GUI and CLI views.}
  \label{fig:case-ops}
\end{figure*}

\paragraph{Spatial and CAD workflow: preserving semantic layers through visual outputs.}
Figure~\ref{fig:case-spa} shows a LibreCAD task. The goal is not only to produce a drawing that looks plausible, but to restore semantic layer assignments and generate release-ready visual evidence. LongHorizon-Harness uses CLI rendering to check that the drawing is materially present, then uses the LibreCAD GUI to inspect and repair layer attributes, and finally validates the print preview and saved screenshots. The auditor checks file size, resolution, uniqueness, and timestamp ordering, ensuring that the submitted screenshots are real captures of distinct states rather than recycled images. This case shows how MEA supports tasks whose success condition spans structured files and visual application state.

\begin{figure*}[t]
  \centering
  \includegraphics[width=\textwidth]{pic/SPA_en.pdf}
  \caption{Spatial/CAD case. The harness combines DXF-level inspection, LibreCAD layer repair, print-preview evidence, and screenshot-integrity checks.}
  \label{fig:case-spa}
\end{figure*}

\paragraph{Design workflow: grounding subjective visual comparison in real rendering state.}
Figure~\ref{fig:case-des} shows a GIMP-based design task involving image upscaling and quality comparison. The executor does not merely produce a claimed result. It opens the actual outputs in GIMP, captures the Levels histogram, compares super-resolution and bicubic results side by side, and uses the real rendering pipeline to generate the evidence. The auditor can then evaluate whether the visual comparison is grounded in the expected files and application state. This case is important because design tasks often contain subjective-looking requirements; LongHorizon-Harness makes them more reliable by converting them into reproducible, inspectable evidence.

\begin{figure*}[t]
  \centering
  \includegraphics[width=\textwidth]{pic/DES_en.pdf}
  \caption{Design case. LongHorizon-Harness grounds visual quality comparison in real GIMP state, including histograms and side-by-side rendered outputs.}
  \label{fig:case-des}
\end{figure*}

\subsection{Terminal-Bench Cases}
\label{app:terminal-additional-cases}

\paragraph{Build and installation: making hidden acceptance criteria explicit.}
Figure~\ref{fig:case-terminal-sqlite} compares baseline and LongHorizon-Harness on a Terminal-Bench task that requires compiling SQLite with gcov instrumentation and making the result available in the shell environment. The baseline reaches partial progress but fails the final reward, illustrating a common long-horizon failure: the agent performs substantial work while missing a blocking acceptance condition. LongHorizon-Harness converts the task into a stable contract covering source provenance, build location, gcov instrumentation, PATH visibility, executability, and preservation of the vendored tarball. Each round audits one subset of this contract before the manager advances. The final delivery is therefore not just a build log, but a verified installation whose binary, PATH entry, gcov symbols, and coverage files are independently checked.

\begin{figure*}[t]
  \centering
  \includegraphics[width=0.92\textwidth]{pic/terminal_en_case2.pdf}
  \caption{Terminal-Bench build case. LongHorizon-Harness turns a multi-condition build task into an explicit acceptance contract and verifies the final SQLite installation with gcov evidence.}
  \label{fig:case-terminal-sqlite}
\end{figure*}

\paragraph{Reverse engineering: auditing hypotheses until they become an independent artifact.}
Figure~\ref{fig:case-terminal-reverse} shows a reverse-engineering task where the agent must produce an independent C program matching a mystery binary's image output under strict constraints. This task is difficult because early observations are hypotheses rather than deliverables: strings, symbols, pixel statistics, and disassembly traces can guide the solution, but none of them alone proves completion. LongHorizon-Harness keeps these observations as audited facts, then uses later rounds to synthesize, compile, compress, and test the final program against the original behavior. The auditor checks not only that the image exists and compiles, but also that the generated output is equivalent when the original binary is removed. This case shows how MEA supports exploratory tasks: uncertain evidence can accumulate without being mistaken for completion, and the final state is accepted only when the independent artifact satisfies the full contract.

\begin{figure*}[t]
  \centering
  \includegraphics[width=0.92\textwidth]{pic/terminal_en_case3.pdf}
  \caption{Terminal-Bench reverse-engineering case. The harness preserves exploratory evidence as audited facts and verifies the final independent C implementation against behavior, compilation, size, and independence constraints.}
  \label{fig:case-terminal-reverse}
\end{figure*}

\subsection{Cross-Case Takeaways}
\label{app:case-study-takeaways}

Across these cases, LongHorizon-Harness improves reliability through a common structure rather than through benchmark-specific heuristics. In document and CAD tasks, it separates visual plausibility from semantic correctness. In web, data-analysis, and DevOps tasks, it forces evidence from different interfaces to agree before marking progress complete. In desktop and design tasks, it preserves recoverable intermediate artifacts when GUI interactions fail. In Terminal-Bench tasks, it converts implicit grading requirements into explicit acceptance contracts and audits them round by round. These patterns support the central claim of the paper: long-horizon agent performance is limited not only by the model's local capability, but also by how the harness represents, verifies, and carries forward task state.



